\section{Discussion}

This study makes two contributions of different character. The first is a \emph{coordinate system}---a log-frequency partition of the EEG spectrum whose boundary positions are empirically optimised, whose scaling ratio is identified by subject-level bootstrap, and whose validity is tested against alternatives with formal model comparison. The second is a \emph{biomarker}---spectral differentiation, a class of within-band individual-differences measures that emerges only when the coordinate system is well-calibrated. The coordinate system is methodological infrastructure; the biomarker is the empirical payload.

These contributions rest on three levels of evidence that should not be conflated. \textbf{Directly observed:} log-frequency geometry, five genuine density troughs, within-band enrichment shapes, age and cognitive and psychopathology associations, test--retest reliability. \textbf{Inferential interpretations:} the alpha mountain as thalamocortical resonance, the $\sim$20~Hz bridge as convergent biophysical time constants, the developmental tuning account of the inverted-U. \textbf{Mechanistic hypotheses to test:} PV+ interneuron maturation as the driver of $\alpha$/$\beta$ trough deepening, GABAergic manipulation predictions, thalamocortical pathology sensitivity, source-localised regional specificity. The IRASA replication provides the sharpest calibration of what is robust: the strongest enrichment structure, the core cognitive association (beta-low differentiation $\times$ reasoning), and the developmental trajectory all survive a fundamentally different aperiodic decomposition method; the broader cognitive map and theta/gamma fine structure do not. The sections below treat each contribution in turn, and Section~3.7 formalises the distinction between robust and provisional claims.

\subsection{Logarithmic Frequency Scaling and the Global Scaling Constant}

The central empirical question of this paper---whether EEG band boundaries follow a logarithmic or linear progression---is answered decisively by the data. Across every dataset and every boundary estimation method, the log-linear model outperforms its arithmetic counterpart by $\Delta$BIC $> 50$, a difference conventionally interpreted as overwhelming evidence. The brain does not divide its spectral territory into equal-bandwidth slices; it divides it into equal-ratio slices, consistent with decades of informal practice in the field and with the theoretical arguments advanced by \textcite{penttonen2003natural}, \textcite{pletzer2010golden}, and \textcite{buzsaki2004neuronal}.

A BIC model comparison across seven candidate ratios initially suggested that the specific scaling constant was ambiguous: $\phisym$, $\sqrt{2}$, and $e - 1$ were mutually indistinguishable ($\Delta$BIC $< 2$). The subject-level bootstrap (Section~\ref{sec:bootstrap}) resolves this ambiguity by shifting the estimand from individual boundary ratios (which are heterogeneous, ranging from 1.36 to 1.89) to the \emph{global} geometric mean---the multiplicative constant that best summarises the five-trough sequence as a whole. At this level, the 95\% CI is [1.609, 1.623], containing $\phisym = 1.618$ while excluding $\sqrt{2}$, $e - 1$, and all other named candidates (all $p < 0.0001$).

This is a strong result but requires careful scoping. It establishes $\phisym$ as the uniquely supported \emph{global coordinate constant} among the candidate theories tested---not as the only possible real number near 1.618. The CI excludes the named competitors but does not exclude, say, 1.615 or 1.620. The claim is: \textbf{of the mathematical constants proposed in the frequency-scaling literature ($e$, $\phisym$, $e-1$, $\sqrt{2}$, $\sqrt{3}$, 2, $2^{1/3}$), only $\phisym$ is compatible with the subject-level evidence.} It does not follow that every adjacent boundary pair is individually $\phisym$-spaced; the local ratios remain irregular and hierarchical. The hierarchy has a coarse scale (three major troughs with ratio $\approx \phisym^2 \approx 2.6$) and a fine scale (five troughs with global ratio $\approx \phisym$), and both are real.

One candidate is excluded with particular force. The natural base $e \approx 2.718$ proposed by \textcite{penttonen2003natural} as the fundamental scaling constant implies inter-band ratios far wider than observed ($\Delta$BIC $> 11$ relative to $\phisym$; excluded by the bootstrap at $p < 0.0001$). The original observation appears to have been based on a selected subset of known band centres rather than systematic peak detection.

\paragraph{Reconciliation with Ursachi (2026).}
\textcite{ursachi2026boundaries} reported that three independent data-driven boundary-detection methods converge on a ratio of $e - 1 \approx 1.718$ rather than $\phisym \approx 1.618$, using $N = 244$ subjects. Our pooled BIC analysis favors $\phisym$ ($-24.74$ vs.\ $-23.06$ for $e - 1$), but this difference is not decisive ($\Delta$BIC $= 1.68$). The two studies' point estimates ($\phisym = 1.618$ vs.\ $e - 1 = 1.718$) differ by only 6\%---well within the variability of per-dataset ratio estimates (log-SD $= 0.26$). Our subject-level bootstrap excludes $e - 1$ at $p < 0.0001$, which is a genuine disagreement with Ursachi's point estimate rather than a methodological nuance. The most likely explanation is that Ursachi's methods detect only the three deepest troughs (geometric mean ratio $\approx 2.5$, biased toward higher values), while our larger pooled sample resolves the two shallower boundaries that bring the global ratio to $\phisym$. Ursachi's power-analysis finding that $N < 30$ systematically favors $\phisym$ does not apply here (our per-dataset $N$s range from 109 to 777), but it appropriately cautions against over-interpreting the specific constant in any single sample. The substantive agreement between the two studies is more important than the constant: both find logarithmic scaling with a ratio in the 1.6--1.7 range, both exclude $e$ and linear spacing, and both confirm that band boundaries are empirical realities rather than analytical conventions.

\subsection{The Enrichment Landscape as an Empirical Contribution}

The within-band enrichment profiles described in Part~II constitute a novel empirical characterization of spectral peak organization that stands independently of the biomarker analyses. Three findings merit emphasis.

First, the two-regime structure---alpha mountain versus ascending ramp---reveals that the canonical frequency bands are not organized by a single mechanism. Alpha's enrichment profile is the precise inverse of all other bands ($r \approx -0.95$ with theta), creating a three-zone octave structure in which the mid-octave belongs exclusively to alpha while the upper octave belongs to everything else. This sign-inversion is consistent with the known thalamocortical origin of the alpha rhythm but has not, to our knowledge, been documented quantitatively at this resolution across multiple datasets. It implies that aggregate enrichment measures---which average across bands---systematically cancel the alpha mountain against the ramps, producing misleadingly small effects. The per-band analysis is not merely more informative; it is necessary.

Second, the cross-boundary typology (cliff, void, bridge, weak) provides a new vocabulary for describing how adjacent bands interact at their shared boundary. The $\sim$20~Hz bridge---the only boundary enriched from both sides---is the single strongest cross-band spectral feature and is paradigm-independent (present in all 9 resting-state datasets). Whether it reflects Fibonacci three-wave resonance \parencite{kramer2023golden,roopun2008period} or convergent biophysical time constants remains open, but its existence constrains any theory of beta-band organization.

Third, the ascending ramp structure (peaks accumulating at the upper edge of each octave in theta, beta-low, and gamma) is a population-level regularity that has not been previously reported. It is not predicted by any existing framework---neither the anti-mode-locking theory \parencite{pletzer2010golden} (which predicts Noble$_1$ enrichment, observed only in alpha) nor the hierarchical coupling framework \parencite{buzsaki2004neuronal} (which predicts boundary enrichment, observed only at the $\sim$20~Hz bridge).

These observations suggest---as a framework for generating testable predictions, not a demonstrated result---that two distinct processes may govern within-band profile shape. The first is \textbf{thalamocortical resonance}: T-type calcium channels in thalamic relay neurons generate intrinsic oscillations in the 8--12~Hz range \parencite{klimesch2018frequency}, acting as a frequency attractor that produces the alpha mountain. Alpha is the only band with a known intrinsic resonator of this type. The second is a \textbf{frequency gradient} in which oscillatory peaks drift toward the upper edge of each octave in the absence of a dedicated resonator---possibly driven by nonlinear harmonic generation, cross-frequency coupling to the next-higher band, or competitive dynamics with the $1/f$ background. The present data cannot distinguish among these candidate gradient mechanisms; specifying the ramp's biophysical basis is an important open problem. The cleanest experimental test would be pharmacological: ethosuximide (a T-type calcium channel blocker) should selectively flatten the alpha mountain without affecting the ascending ramps of other bands.

\paragraph{The trough depth hierarchy and its developmental trajectory.}
The five band boundaries differ not only in their frequency positions but in their depth: two troughs show greater than 60\% peak depletion (5.1~Hz, 70.4\%; 13.4~Hz, 61.7\%) while three are substantially shallower (7.8~Hz, 8.7\%; 25.3~Hz, 11.6\%; 35.0~Hz, 32.2\%; Table~\ref{tab:aperiodic_null}). This bimodal depth pattern is not predicted by the coordinate-geometric account developed in Section~3.3---boundaries that serve only as measurement partitions should not vary in depth. Moreover, when measured across age bins within the HBN dataset (supplementary trough-depth-by-age analysis), the five troughs fall into three developmental maturation classes at the youngest ages tested (5--7 years): \emph{over-mature} ($\delta$/$\theta$ at 233\% of adult depth, $\beta_L$/$\beta_H$ at 297\%), \emph{immature} ($\alpha$/$\beta$ at 40\% of adult depth, $\theta$/$\alpha$ at 31\%), and \emph{absent} ($\beta_H$/$\gamma$ at 3\%). Both the depth hierarchy and the maturation hierarchy require explanation beyond the coordinate framework.

The $\alpha$/$\beta$ trough at 13.4~Hz provides the clearest developmental signal. Within HBN alone (ages 5--21, $N = 927$), trough depth increases monotonically with age ($\rho = +0.82$, $p = 0.023$ across 7 age bins, each pooling 37--220 subjects), from approximately 20\% depletion at age 8 to 50\% at age 19; the per-subject continuous correlation confirms the direction ($\rho = +0.31$, $p < 0.0001$, $N = 923$), with the smaller magnitude reflecting the noise inherent in individual trough-depth estimates from $\sim$2{,}000 peaks per subject. In the Dortmund adult data (ages 20--70, $N = 608$), the same trough is flat at $\sim$55--65\% depletion ($\rho \approx 0$, $p = 0.99$), confirming that it reaches a stable plateau by the early twenties. The transition between datasets is smooth, ruling out a cross-cohort artefact. This developmental profile---immature in childhood, deepening through adolescence, stable from early adulthood---matches the known timeline of parvalbumin-positive (PV+) fast-spiking interneuron maturation, in which perineuronal net consolidation completes in the late teens to early twenties in sensory cortex and somewhat later in prefrontal cortex \parencite{hensch2005critical}. If the $\alpha$/$\beta$ trough marks the lower frequency limit of PV+ perisomatic inhibition, then its developmental deepening reflects the progressive sharpening of that frequency floor as inhibitory circuits mature. The $\alpha$/$\beta$ trough also forms a weak positive covariance cluster with the $\beta_H$/$\gamma$ trough across individuals ($\rho = +0.37$), consistent with shared PV+ fast-spiking contribution across the 13--40~Hz range, but is essentially independent of the $\delta$/$\theta$ trough ($\rho = +0.08$, NS within Dortmund). In HBN, the $\alpha$/$\beta$ trough shows the psychopathology dissociation predicted by an inhibitory account: higher externalising symptom scores are associated with a shallower trough ($\rho = +0.146$, FDR $q < 0.05$), while higher internalising scores are associated with a deeper trough ($\rho = -0.116$, FDR $q < 0.05$), replicating the within-band enrichment dissociation (Section~\ref{sec:psychopathology}) at the boundary level.

The $\delta$/$\theta$ trough at 5.1~Hz tells a different story. Despite being the deepest trough in the pooled data, it is \emph{over-deep} in young children (233\% of adult depth at age 6) and \emph{decreases} toward adult levels during development ($\rho = -0.67$, $p = 0.013$ within HBN), following a trajectory anti-correlated with the $\alpha$/$\beta$ trough ($r = -0.58$, $p = 0.037$). This pattern is inconsistent with inhibitory maturation of any kind and most parsimoniously reflects the developmental recession of delta-band oscillatory power: young children generate dominant slow-wave activity that creates sharp spectral peaks below 5~Hz, deepening the $\delta$/$\theta$ gap; as delta power diminishes with age, the trough fills in. The developmental data therefore suggest two mechanistically distinct classes of trough: \emph{inhibitory-boundary troughs} ($\alpha$/$\beta$, possibly $\theta$/$\alpha$) that deepen during development as inhibitory circuits mature, and \emph{generator-dominated troughs} ($\delta$/$\theta$, $\beta_L$/$\beta_H$) whose depth tracks the waxing and waning of the oscillatory generators flanking them.

These developmental and covariance findings extend the pharmacological predictions above. Benzodiazepines (GABA$_\text{A}$ positive allosteric modulators) should deepen the 13.4~Hz trough---where PV+ inhibition is proposed as the boundary mechanism---and shift it downward in frequency, while leaving the 5.1~Hz trough unaffected, since that boundary is generator-dominated rather than GABAergic. Tiagabine (a GABA reuptake inhibitor) should broaden the $\alpha$/$\beta$ trough without affecting the $\delta$/$\theta$ trough. The differential response of the two deep troughs to GABAergic agents would directly test the two-class mechanism model. These predictions, like the two-mechanism and ethosuximide predictions above, are post hoc and lack the specificity to exclude alternative developmental accounts. In particular, myelination-dependent increases in conduction velocity could produce similar trough-deepening profiles without invoking PV+ circuits specifically: faster conduction sharpens the temporal precision of oscillatory networks, which could concentrate spectral peaks and deepen inter-band troughs on a timeline that overlaps with PV+ maturation. The pharmacological predictions discriminate between these accounts---GABAergic agents should affect the $\alpha$/$\beta$ trough if it is inhibition-dependent, but not if it is myelination-dependent---and their value lies in their joint falsifiability. These supplementary trough-depth analyses are summarised in Figure~\ref{fig:trough_analyses}.

\begin{figure}[htbp]
\centering
\includegraphics[width=\textwidth]{fig_trough_analyses.png}
\caption{\textbf{Trough depth analyses.} (A)~$\alpha$/$\beta$ trough depletion across the lifespan. Within HBN (ages 5--21), depletion increases monotonically ($\rho = +0.82$, $p = 0.023$); within Dortmund (ages 20--70), it is stable ($\rho \approx 0$, $p = 0.99$). Shaded bands: 95\% bootstrap CIs from 500 subject-level resamples per age bin. (B)~Differential maturation: the $\delta$/$\theta$ trough is over-deep in children ($>$200\% of adult depth) and regresses, while the $\alpha$/$\beta$ trough is immature ($\sim$40\%) and deepens on the PV+ interneuron maturation timeline. Dashed line: adult reference (Dortmund 25--55). (C)~Per-subject trough depth correlation matrix ($N = 1{,}738$). The five troughs are largely independent (mean off-diagonal $\rho = 0.053$); the $\alpha$/$\beta$--$\beta_H$/$\gamma$ cluster ($\rho = 0.37$) is the only substantial positive correlation. (D)~Psychopathology $\times$ trough depth in HBN ($N = 906$). Externalising and internalising symptoms push the $\alpha$/$\beta$ trough in opposite directions ($\rho = +0.146$ and $-0.116$, both FDR $q < 0.05$), consistent with GABA deficit and enhancement models respectively. ${}^{***}p < 0.001$.}
\label{fig:trough_analyses}
\end{figure}

\subsection{Boundaries Matter; Landmarks Do Not}

A conceptual clarification is warranted by the within-band profile results. Prior formulations of the $\phisym$-lattice framework \parencite{pletzer2010golden,kramer2023golden} have emphasized that noble-number frequency ratios suppress mode-locking between oscillators, suggesting that specific frequencies carry dynamical significance as anti-resonance points. The present data do not support this interpretation at the within-band level. Within-band frequency profiles are smooth and monotone---ascending ramps in most bands, a single mountain in alpha---with no evidence of fine structure, enrichment peaks, or depletion troughs at the positions predicted by noble-number arithmetic. Landmark frequencies, in this sense, do not matter.

What does matter is where the boundaries between bands are drawn. Log-spaced boundaries produce the simplest, most interpretable within-band spectral profiles. Misaligned boundaries---whether from using linear spacing or from individual $\fzero$ mismatch---import spectral content from adjacent bands, fragmenting what would otherwise be a clean profile and generating spurious enrichment and depletion patterns. The value of $\phisym$-scaling is coordinate-geometric rather than dynamical: it provides a partition of the frequency axis that minimizes cross-band leakage and thereby maximizes the discriminability of within-band oscillatory organization. A reviewer might note that the third-octave system ($r = 1.26$) achieves higher cross-dataset consistency ($r = 0.87$ vs.\ $0.79$ for the $\phisym$-lattice) in the boundary sweep. We optimize for profile simplicity rather than raw consistency because simplicity---the ability to describe within-band enrichment with a low-order polynomial---is what makes individual-differences metrics interpretable: a monotonic ramp has a meaningful slope, a mountain has a meaningful height, and both can be summarized by single numbers that predict cognition. High cross-dataset consistency of a \emph{complex} profile would be less useful, because the derived metrics would require more parameters and be harder to validate. The $\phisym$-lattice produces the simplest profiles ($R^2 = 0.968$) precisely because its boundaries fall where the spectral structure naturally transitions, yielding clean ramps and mountains rather than multi-modal shapes.

This reframing---from ``the brain computes with noble numbers'' to ``$\phisym$-scaling provides measurement-optimal boundaries for quantifying spectral differentiation''---is empirically more defensible and practically more consequential.

\subsection{Spectral Differentiation as a Novel Biomarker Class}

The spectral differentiation index introduced here occupies a distinct niche in the landscape of EEG-derived measures. It is not band power, which quantifies total energy without regard to internal structure. It is not the individual alpha frequency \parencite{klimesch1999eeg,grandy2013iaf,corcoran2018iaf}, which identifies a single peak location. It is not the aperiodic slope \parencite{donoghue2020fooof,voytek2015aperiodic}, which characterizes the background $1/f$ process. And it is not connectivity, which concerns relationships between spatially separated signals. Spectral differentiation measures something orthogonal to all of these: how organized the within-band peak distribution is, treating the histogram of individual-epoch spectral peaks as the object of inference.

A band with high spectral differentiation is one in which oscillatory epochs cluster tightly around a characteristic frequency, with few epochs straying into neighboring sub-bands. A band with low differentiation is one in which oscillatory activity is diffuse, irregular, or absent as a coherent phenomenon. Why should this matter for cognition? We suggest---as a hypothesis, not a demonstrated mechanism---that tighter frequency tuning supports more efficient neural communication through at least two routes. First, narrower frequency spread enables faster phase-locking between oscillating populations: two circuits operating at 19.5 $\pm$ 0.3 Hz can synchronize within a few cycles, while circuits operating at 19.5 $\pm$ 2.0 Hz require substantially longer to achieve coherence. Second, precise frequency tuning sharpens cross-frequency coupling: a 20 Hz signal couples to its 40 Hz harmonic more reliably than a smeared 17--23 Hz signal couples to anything in the gamma range. Both routes predict that high differentiation supports more efficient information transfer, consistent with the observed correlations with reasoning and executive function. This account is consistent with the ``communication through coherence'' framework \parencite{buzsaki2004neuronal} but remains speculative until tested with concurrent oscillatory and behavioral measures.

The cognitive correlate results underscore the biomarker potential. The core cognitive finding is that beta-low center depletion predicts reasoning ability ($\rho = -0.27$), survives both FDR correction and age partialing, and replicates under IRASA ($\rho = -0.29$, FDR $q < 0.02$). This cross-method replication anchors the cognitive story: it is not a FOOOF artefact. Beyond this anchor, an exploratory cognitive map generated by the framework identifies 31 FDR-surviving associations across four frequency bands and four cognitive domains in the LEMON dataset ($N = 203$). Earlier extraction pipelines with overly restrictive peak caps found zero cognitive FDR survivors in the same dataset---the 31 survivors emerge from correcting a systematic artefact, not from analytical degrees of freedom. Critically, cognitive associations survive IAF-partialing with only 23\% attenuation (vs.\ 47\% from age), confirming that spectral differentiation is not reducible to boundary misalignment. These effects are modest (peak $|\rho| \approx 0.27$, $R^2 \approx 7\%$ raw, $\approx 2\%$ unique after age), comparable to the strongest IAF--cognition effects in the literature ($|\rho| \approx 0.25$--$0.35$; \textcite{grandy2013iaf}). The broader 31-association map is FOOOF-specific and single-dataset; it should be regarded as a discovery-layer context for future replication, not a settled atlas.

Equally important are the informative nulls. Spectral differentiation is uninformative about personality traits (0/11,970 FDR in LEMON), medical and metabolic variables (0 FDR across $\sim$2,500 tests), and handedness (0 FDR in HBN and Dortmund). The contrast between 31 cognitive FDR survivors and zero personality survivors, in the same subjects with the same features, argues against the construct being a generic nuisance correlate or a signal artefact. It is selectively predictive of cognitive performance, not globally correlated with individual differences.

\subsection{An Inhibitory Circuit Hypothesis}
\label{sec:inhibitory}

The phase-locking and cross-frequency-coupling routes described above address the proximal mechanism by which differentiation supports cognition---tighter tuning enables faster synchronisation. A complementary question is what produces the tuning itself. This section consolidates the mechanistic evidence from the trough-depth analyses (Section~3.2), the psychopathology dissociation (Section~\ref{sec:psychopathology}), and the developmental trajectory (Sections~\ref{sec:development}--\ref{sec:lifespan}) into a unified hypothesis-generating framework. We suggest---as a hypothesis with testable pharmacological consequences, not a demonstrated result---that spectral differentiation is primarily a readout of PV+ fast-spiking interneuron circuit integrity, expressed most clearly at the $\alpha$/$\beta$ boundary. A brain with high spectral differentiation has sharp inhibitory boundaries that confine oscillatory peaks to well-defined spectral territories; a brain with low differentiation has weak or immature inhibitory boundaries that allow spectral energy to smear across the frequency axis.

This interpretation connects four otherwise disparate findings. First, the externalising--internalising psychopathology dissociation (Section~\ref{sec:psychopathology}): externalising symptoms (associated with GABAergic deficits in ADHD; \textcite{edden2012gaba}) are associated with flatter within-band profiles \emph{and} a shallower $\alpha$/$\beta$ trough, while internalising symptoms are associated with more concentrated upper-alpha enrichment \emph{and} a deeper $\alpha$/$\beta$ trough---the direction expected if anxious hypervigilance reflects enhanced, not diminished, inhibitory tone. Second, the personality null (Section~\ref{sec:nulls}): personality dimensions lack established GABAergic substrates, so a readout of inhibitory circuit integrity should be silent for personality, which it is (0/11,970 FDR). Third, the cognitive results operate at the within-band level, not the boundary level: supplementary analyses found only 1/40 FDR-surviving associations between per-subject trough depth and cognitive scores in LEMON, versus 31 for within-band enrichment. This distinction is expected if the cognitively relevant output of inhibitory circuits is peak \emph{organisation}---how precisely oscillatory activity is tuned within each band---rather than gap \emph{depth} per se. Fourth, the developmental trajectory of the $\alpha$/$\beta$ trough (Section~3.2) directly tracks PV+ maturation.

These predictions are jointly testable. Table~\ref{tab:pharmacological} enumerates three pharmacological agents, each targeting a different component of the proposed mechanism. The pattern of selective effects on the $\alpha$/$\beta$ trough (predicted to respond to GABAergic agents) versus the $\delta$/$\theta$ trough (predicted to be unaffected, since it is generator-dominated) would directly discriminate the inhibitory model from alternatives. A complementary correlational test: individual differences in spectral differentiation should correlate with cortical GABA concentration measured by magnetic resonance spectroscopy (MRS); no MRS-EEG study has yet examined within-band peak distributions. The inhibitory interpretation is not uniquely determined by the present data---variation in excitatory mechanisms (e.g., NMDA receptor kinetics) could produce similar patterns---but the GABAergic account is distinguished by its parsimony (a single mechanism class explains both the trough depth hierarchy and the psychopathology dissociation) and its experimental accessibility.

\begin{table}[h]
\centering
\caption{Predicted pharmacological dissociations under the GABAergic inhibitory model. Each agent targets a different component of the proposed mechanism; the pattern of selective effects on the $\alpha$/$\beta$ versus $\delta$/$\theta$ trough distinguishes the inhibitory model from alternatives. None of these predictions has been tested.}
\label{tab:pharmacological}
\begin{tabular}{l l l l l}
\toprule
Agent & Target & $\alpha$/$\beta$ trough & $\delta$/$\theta$ trough & Distinguishes \\
\midrule
Ethosuximide & T-type Ca$^{2+}$ & No change & No change & Null on both = inhibitory model \\
Benzodiazepine & GABA$_\text{A}$ PAM & Deepens, shifts down & No change & $\alpha$/$\beta$-only = PV+ model \\
Tiagabine & GABA reuptake & Broadens & No change & $\alpha$/$\beta$-only = two-class model \\
\bottomrule
\end{tabular}
\end{table}

\subsection{Spectral Differentiation and the Inverted-U of Cognitive Aging}

If spectral differentiation tracks the inverted-U of neural differentiation across the lifespan \parencite{li2004dedifferentiation}, it would provide the first frequency-domain EEG measure with a principled connection to this framework. The de-differentiation hypothesis holds that neural representations become more distinct and specialized through childhood and early adulthood, then progressively lose specificity in aging, with this loss mediating much of the cognitive decline observed in older adults.

The cross-sectional HBN developmental data are consistent with the ascending limb: spectral differentiation increases from childhood through adolescence, attributable to the progressive maturation of myelination, synaptic pruning, and inhibitory interneuron populations that together sharpen the frequency tuning of cortical circuits. Of these processes, the PV+ inhibitory interneuron trajectory is the most temporally specific. The $\alpha$/$\beta$ trough reaches 50\% of its adult depth at age 12 and 75--90\% by age 17 (Section~3.2), coinciding with perineuronal net consolidation that stabilises PV+ fast-spiking function \parencite{hensch2005critical}. The distinction between \emph{immature} troughs ($\alpha$/$\beta$, $\theta$/$\alpha$) that deepen during development and \emph{over-mature} troughs ($\delta$/$\theta$, $\beta_L$/$\beta_H$) that regress toward adult levels (Section~3.2) suggests that the ascending limb of the inverted-U has two components: progressive sharpening of inhibitory boundaries and simultaneous recession of generator-dominated boundaries that were over-deep in childhood. LEMON adult data then show the expected positive relationship between differentiation and cognitive performance. The Dortmund aging data show the descending limb: the same features that increase during development decrease during aging, with 28 of 31 cross-dataset features showing opposite signs. At the trough level, the $\theta$/$\alpha$ boundary shallows significantly during aging ($\rho = -0.81$, $p = 0.005$ within Dortmund), and the $\beta_L$/$\beta_H$ trough inverts from depletion to enrichment ($\rho = -0.92$, $p = 0.0002$), suggesting progressive degradation of boundary structure in the second half of life---consistent with the known vulnerability of PV+ interneurons to oxidative stress and aging-related loss.

If spectral differentiation is the oscillatory substrate of cognitive de-differentiation, then aging populations should show systematically lower differentiation scores, and the magnitude of this decline should predict the degree of cognitive impairment. This prediction is testable in existing datasets such as TDBRAIN or Cambridge Centre for Ageing and Neuroscience (Cam-CAN). However, as noted in Section~\ref{sec:lifespan}, the inverted-U trajectory is driven entirely by the cross-cohort splice between HBN and Dortmund rather than by a within-cohort turning point, and longitudinal data spanning childhood through middle age would be needed to confirm whether the vertex at $\sim$33--35 years is genuine or an artefact of the cross-sectional design.

\subsection{Corrections to Prior Work}

Scientific progress occasionally requires explicit correction of one's own prior claims, and this paper presents several. \textcite{lacy2026frontiers} reported two findings as universal features of the $\phisym$-lattice framework: aggregate boundary depletion (spectral troughs co-locating with band boundaries across all bands) and Noble$_1$ enrichment (peaks clustering near the first noble-number landmark within each band). Both claims were supported by statistically significant results in that analysis.

The present v3 analysis demonstrates that both findings were artifacts of the extraction pipeline. The aggregate boundary depletion was driven by a $\fzero$ mismatch that wrapped peaks into incorrect Voronoi bins, most consequentially reversing the beta-low boundary from $+101\%$ to $-33\%$. The Noble$_1$ enrichment is real but alpha-specific---it does not generalize to theta, beta-low, beta-high, or gamma (all show Noble$_1$ $\approx 0\%$). The aggregate effect reported in the Frontiers paper was an average over five bands in which one band (alpha) showed strong enrichment and four showed null, producing a misleadingly positive aggregate.

Earlier extraction pipelines with overly restrictive peak caps produced zero cognitive FDR survivors and negative test--retest ICCs ($\text{ICC} < 0$) for spectral differentiation. The 31 FDR survivors and strongly positive ICCs ($0.42$--$0.75$) reported here emerge from the same datasets with corrected extraction parameters, demonstrating that the prior null results were methodological artifacts rather than genuine absence of signal.

The framework introduced in \textcite{lacy2026frontiers} is confirmed as a measurement framework by these results; the specific empirical claims it advanced required correction. Three methodological corrections substantially change the enrichment landscape and reveal previously hidden cognitive, developmental, and clinical associations.

\subsection{Robust and Provisional Claims}

For clarity, we distinguish claims that are strongly supported by the present evidence from those that remain provisional. \textbf{Robust claims}: logarithmic frequency scaling ($\Delta$BIC $> 50$); five real density troughs confirmed by the aperiodic null ($p < 0.0001$) and recovered in 100\% of subject-level bootstraps; $\phisym$ as the best-supported named global boundary constant tested (bootstrap 95\% CI [1.609, 1.623]); the $\phisym$-lattice as a measurement-optimal coordinate system for profile simplicity ($R^2 = 0.968$; boundaries within 2\% of locally optimal positions); the alpha mountain and upper-octave depletion cascade (11/13 cross-dataset consistency under FOOOF, replicated under IRASA); the beta-high descending ramp (method-independent, 13/13 cross-dataset consistency under IRASA); the top cognitive association---beta-low center depletion as a predictor of reasoning---replicating across methods (FDR-significant under both FOOOF and IRASA), though the breadth of the cognitive map (31 FDR survivors) is FOOOF-specific; the developmental inverted-U trajectory (68 FDR survivors under both methods); 5-year test--retest reliability of beta-low spectral differentiation (ICC $= 0.75$). \textbf{Provisional claims}: exact $\phisym$ ontology (the CI contains $\phisym$ but also nearby values; local boundary ratios are heterogeneous); theta enrichment geometry (method-dependent, inverted under IRASA); gamma within-band architecture (weak cross-dataset consistency, scalp EEG contamination); the specific cognitive effect sizes (single-dataset, modest $\rho$); the vertex of the inverted-U at 33--35 years (cross-sectional splice); the PV+ inhibitory interpretation of the $\alpha$/$\beta$ trough and spectral differentiation (developmental and psychopathological correlations consistent with PV+ maturation, but post hoc and without direct neurochemical evidence; pharmacological and MRS predictions untested); and the two-class trough mechanism model (inhibitory-boundary versus generator-dominated; based on developmental trajectory classification, not direct mechanism measurement).

\subsection{Limitations}
\label{sec:limitations}

Several limitations constrain the conclusions drawn here. First, the cognitive correlate analyses rely on a single dataset (LEMON, $N = 203$) and require independent replication before the 31 FDR-surviving associations can be regarded as established. At $N = 203$, a two-tailed Spearman correlation has approximately 80\% power to detect $|\rho| \geq 0.20$ at $\alpha = 0.05$ for a single test; however, FDR correction across 720 simultaneous tests raises the effective detection threshold, meaning that true associations weaker than the observed $|\rho| \approx 0.25$ are likely missed. The 31 FDR survivors should therefore be regarded as a lower bound on the number of genuine associations rather than an exhaustive catalogue. Sample size, demographic homogeneity, and task battery all limit generalizability.

Second, the band boundaries employed throughout are defined at the population level. Individual $\fzero$ values span 8--13 Hz in healthy adults \parencite{klimesch1999eeg,grandy2013iaf}, sufficient to cause boundary misalignment. The IAF-partialing analysis (Section~\ref{sec:cognitive}) demonstrates that this confound is not the primary driver of the cognitive results: IAF-partialing attenuates correlations by only 23\% (and adds only 4 percentage points beyond age alone, reflecting IAF-age collinearity), with 31/46 strong features surviving. Nevertheless, individual-specific coordinate systems anchored to each subject's $\fzero$ remain the principled approach and should improve sensitivity, particularly for beta-low features where the alpha mountain's flank determines ramp position.

Third, the subject-level bootstrap discriminates among the named candidate constants ($\phisym$, $e-1$, $\sqrt{2}$, etc.) by showing that only $\phisym$ lies within the 95\% CI for the global inter-trough ratio. However, it does not establish that $\phisym$ is the exact value---the CI [1.609, 1.623] also contains nearby non-named values---nor does it imply exact local $\phisym$-spacing at every boundary, since individual boundary ratios remain heterogeneous (1.36--1.89). The bootstrap resolves the named-constant question but not the deeper question of whether the ratio is exactly $\phisym$ or merely a biophysical constant that happens to fall near 1.618.

Fourth, the developmental data from HBN are cross-sectional. Longitudinal tracking of spectral differentiation from childhood through adulthood would be required to establish the causal direction of the differentiation-cognition relationship and to test the inverted-U prediction directly.

Fifth, all primary analyses rely on FOOOF/specparam for aperiodic-periodic separation. The IRASA replication (Section~\ref{sec:irasa}) demonstrates selective method-independence: beta-high enrichment, the top cognitive association, and the developmental trajectory replicate under a fundamentally different non-parametric approach, while theta and gamma remain method-sensitive. The broader cognitive map (31 FDR survivors) is FOOOF-specific; a peak-yield subsample test (Section~\ref{sec:irasa}) attributes this primarily to IRASA's lower per-subject peak yield in the bands carrying the strongest cognitive signal. The strongest features are not artefacts of the parametric approach; the broader map requires either FOOOF's yield or larger samples.

\subsection{Future Directions}

The most immediate priority is the construction of individual-specific frequency coordinate systems. If each subject's band boundaries are shifted proportionally to their $\fzero$---so that the theta-alpha boundary falls at $\fzero / \phisym$ regardless of whether $\fzero = 9$ Hz or $\fzero = 11$ Hz---the boundary misalignment artifacts are eliminated by construction. This approach, which we term IAF-anchored spectral differentiation, should be the default in future work \parencite{corcoran2018iaf,klimesch1999eeg}.

Source-space analysis represents a second priority. MEG, ECoG, and high-density EEG with source modeling would allow spectral differentiation to be applied to specific cortical and subcortical structures, testing whether cognitive correlates show regional specificity and resolving the gamma contamination concern.

Clinical translation offers a third direction. Neurodegenerative diseases are associated with progressive disruption of thalamocortical circuits \parencite{klimesch2018frequency}, and the two-mechanism model predicts that alpha mountain flattening should be among the earliest spectral signatures of thalamic pathology.

Pharmacological and neurochemical validation of the inhibitory model represents a fourth priority. The predictions in Table~\ref{tab:pharmacological} specify a crossover design in which ethosuximide, a benzodiazepine, and tiagabine are administered to healthy volunteers, with pre- and post-drug spectral differentiation and trough depth computed from resting EEG. The differential prediction---GABAergic agents affecting the $\alpha$/$\beta$ trough but not the $\delta$/$\theta$ trough---directly tests the two-class mechanism model. An MRS-EEG study correlating individual cortical GABA concentration with spectral differentiation would provide complementary correlational evidence and represents the most immediately feasible test of the framework, given that MRS-GABA protocols are established \parencite{edden2012gaba} and could be applied to existing multimodal datasets such as Cam-CAN.

Finally, the neurofeedback literature has established that oscillatory power can be upregulated through closed-loop training, but has not examined whether spectral differentiation---the organization of oscillatory activity, as distinct from its magnitude---can be similarly targeted. If the cognitive benefits of neurofeedback derive from improved spectral organization rather than increased power, then differentiation-targeted protocols should outperform power-targeted ones.
